%! Polar Function Graphs — Unit 3, Lesson 3.14 — Homework (Answer Key)
\documentclass[10pt]{article}
\usepackage{precalculus-article}
\usepackage{precalculus-key}   % pulls in -boxes; do NOT also load -boxes
\newcommand{\TallMath}[1]{$\displaystyle #1\rule[-1.4em]{0pt}{3.2em}$}

\begin{document}

\pageheader{Unit 3, Lesson 3.14}{Homework --- Answer Key}
\namedateperiod

\vspace{0.10in}

\begin{notesbox}{Problems 1--6 --- Key}

\textbf{1.}\quad Table for $r = 1 - \sin\theta$.

\vspace{4pt}
\renewcommand{\arraystretch}{1.55}
\begin{tabular}{|c|c|c|c|c|c|c|c|}
\hline
$\theta$ & $0$ & $\dfrac{\pi}{6}$ & $\dfrac{\pi}{2}$ & $\pi$ &
           $\dfrac{3\pi}{2}$ & $\dfrac{5\pi}{4}$ & $2\pi$ \\
\hline
$r = 1 - \sin\theta$ & $1$ & $0.5$ & $0$ & $1$ & $2$ & $1.71$ & $1$ \\
\hline
\end{tabular}

\vspace{6pt}
(a)\quad Maximum $r$: \ans{2}, at $\theta = $ \ans{$3\pi/2$}.

\vspace{4pt}
(b)\quad $r = 0$ at $\theta = $ \ans{$\pi/2$}. The curve passes through the \ans{pole (origin)}.

\vspace{4pt}
(c)\quad On $\theta \in \left[0, \frac{\pi}{2}\right]$, $r$ is \ans{decreasing} (from 1 to 0).

\vspace{0.14in}
\textbf{2.}\quad Rectangular graph of $r = \cos(3\theta)$ for $\theta \in [0, \pi]$.

\vspace{4pt}
\begin{center}
\begin{tikzpicture}[x=1.55cm, y=1.1cm]
  \draw[linegray, very thin] (-0.2,-1.3) grid[xstep=0.25, ystep=0.5] (3.3,1.3);
  \draw[->] (-0.2,0) -- (3.45,0) node[right]{\small $\theta$};
  \draw[->] (0,-1.35) -- (0,1.4) node[above]{\small $r$};
  \foreach \x/\lbl in {0.524/$\frac{\pi}{6}$, 1.047/$\frac{\pi}{3}$,
                        1.571/$\frac{\pi}{2}$, 2.094/$\frac{2\pi}{3}$,
                        2.618/$\frac{5\pi}{6}$, 3.142/$\pi$} {
    \draw (\x,0.07) -- (\x,-0.07);
    \node[below, font=\scriptsize] at (\x,-0.07) {\lbl};
  }
  \foreach \y/\lbl in {-1/{$-1$}, 1/{$1$}} {
    \draw (0.07,\y) -- (-0.07,\y);
    \node[left, font=\scriptsize] at (-0.07,\y) {\lbl};
  }
  \draw[navy, thick, domain=0:3.14, samples=200]
    plot (\x, {cos(3*deg(\x))});
  \node[right, font=\small, navy] at (3.2, 0.5) {$r = \cos(3\theta)$};
\end{tikzpicture}
\end{center}

\vspace{4pt}
(a)\quad $r = 0$ at: \ans{$\theta = \pi/6,\ \pi/2,\ 5\pi/6$} (within $[0,\pi]$; also at $\pi/3 \cdot$ multiples of $\pi/2$, specifically $\pi/6$, $\pi/2$, $5\pi/6$).

\vspace{4pt}
(b)\quad \ans{3} complete arcs in $[0, \pi]$.

\vspace{4pt}
(c)\quad The polar curve $r = \cos(3\theta)$ has \ans{3} petals.

\vspace{0.14in}
\textbf{3.}\quad Domain restriction for $r = \sin(2\theta)$.

\vspace{4pt}
(a)\quad \ans{1} petal traced for $\theta \in [0, \pi/2]$.

\vspace{4pt}
(b)\quad \ans{For $\theta \in [0, \pi]$, the rectangular graph completes two full arcs (one positive, one negative). The polar curve traces 2 petals: one in the first quadrant (from the positive arc) and one in the third quadrant (from the negative arc, since negative $r$ plots in the opposite direction).}

\vspace{0.14in}
\textbf{4.}\quad AP-Style MC: \ans{(B) 4}.

\begin{teachernote}
Maximum of $r = 2 + 2\sin\theta$ occurs when $\sin\theta = 1$ (at $\theta = \pi/2$):
$r_{\max} = 2 + 2(1) = 4$.
\end{teachernote}

\vspace{0.14in}
\textbf{5.}\quad AP-Style MC: \ans{(C) 8}.

\begin{teachernote}
$n = 4$ (even), so petal count $= 2n = 2(4) = 8$.
\end{teachernote}

\vspace{0.14in}
\textbf{6.}\quad $r = 3$ (constant polar function).

\vspace{4pt}
(a)\quad The radius is always \ans{3}.

\vspace{4pt}
(b)\quad The curve traces a \ans{circle of radius 3 centered at the origin}.

\vspace{4pt}
(c)\quad Domain needed: \ans{$\theta \in [0, 2\pi]$} (one full revolution traces it once).

\end{notesbox}

\vspace{0.08in}

\begin{tcolorbox}[colback=navylight, colframe=navy, arc=2mm,
    left=3mm, right=3mm, top=2mm, bottom=2mm]
\small\textbf{Preview --- Lesson 3.15: Rates of Change of Polar Functions}\\
You can now read a polar function from tables and rectangular graphs. In Lesson 3.15
we'll ask: how fast is the radius $r$ changing as the angle $\theta$ increases?
Come ready to connect the slope of the rectangular graph to the behavior of the
polar curve.
\end{tcolorbox}

\end{document}
